% !TeX root = ../../Book1.tex
\chapter{余面积公式}
\label{b1:ch:18}
% 本章摘要（不计入编号）
\begin{chapterabstract}
把高维区域按映射值分层，得到的不是离散重数，而是每条纤维自身的低维面积。本章从水平集的线性化和正规坐标出发，完整证明 Lipschitz 映射的余面积公式。零测集的切片、临界部分的贡献和切片积分的可测性分别处理；应用中再讨论球壳积分、水平集面积的平均控制，以及分布密度公式为什么不能无条件除以梯度。
\end{chapterabstract}

% 本章目标（不计入编号）
\begin{learninggoals}
\begin{itemize}
\item\label{b1:item:ch18-goal-level}理解正规水平集的切空间与法空间，区分光滑局部坐标和一般 Lipschitz 水平集。
\item\label{b1:item:ch18-goal-jacobian}能比较面积与余面积两种 Jacobi 因子的维数及矩阵形式，计算切片坐标中的补偿因子。
\item\label{b1:item:ch18-goal-proof}掌握余面积证明中零测切片、临界切片和满秩分片的分工，并能使用非负及可积加权版本。
\item\label{b1:item:ch18-goal-density}能从水平集计算积分与分布，保留临界质量，并分清超水平集体积与等值集面积。
\end{itemize}
\end{learninggoals}

本章固定定义域维数 \(m\) 不小于目标维数 \(n\)。公式把 \(m\) 维体积拆成 \((m-n)\) 维纤维面积与 对 \(n\) 维参数的积分；当两维相等，纤维上的计数测度使它回到前章的面积公式。读者应先熟悉 Tonelli 定理、第十三章的 Besicovitch 覆盖及第十七章的单射面积公式。第\ref{b1:sec:1801}节把光滑正规坐标所需的局部求逆补证出来；一般 Lipschitz 证明则不假定这些坐标存在于整块开邻域中，而在可数个双 Lipschitz 片上进行。

Simon 的《Lectures on Geometric Measure Theory》第10节\cite{Simon1983Lectures}可供核对线性余面积及水平集的几何表述；它把部分逼近论证交给外部文献，本书则利用前章工具补齐 Lipschitz 主线。Fremlin 的《Measure Theory》第265G节\cite{Fremlin2016Measure}提供球壳积分的参数化处理，便于与本章的径向特例比较。一般余面积及其后续几何应用可参阅 Federer 的《Geometric Measure Theory》\cite{Federer1996Geometric}；这里推荐的是进一步系统阅读的方向，不把未实际读取的原书证明冒作本章现成依据。

% !TeX root = ../../Book1.tex
\section{水平集}
\label{b1:sec:1801}

% TODO: 撰写本节正文。

\subsection{Lipschitz 函数的水平集}
\label{b1:subsec:180101}

待填充。

\subsection{法方向}
\label{b1:subsec:180102}

待填充。

\subsection{切空间}
\label{b1:subsec:180103}

待填充。

\subsection{余面积 Jacobi 行列式}
\label{b1:subsec:180104}

待填充。

% !TeX root = ../../Book1.tex
\section{余面积公式}
\label{b1:sec:1802}

面积公式沿参数方向把低维区域映入高维空间；\term[yumianji-gongshi]{余面积公式}[coarea formula]则把高维区域按映射值分层。正交投影时，这就是 Fubini 定理。一般 Lipschitz 映射的水平集不再平行，补偿因子也由常数变成余面积 Jacobi 行列式。
线性情形与几何表述可对照 Simon 的《Lectures on Geometric Measure Theory》\cite[第 10 节，第 53--57 页]{Simon1983Lectures}；本节从零测集的切片、满秩坐标分片和前章面积公式出发，证明 Lipschitz 版本。

% 实际读取：本地Simon1983Lectures §10印刷pp.53--57全文。
% 该来源的10.3只给C1版本并把逼近论证交外部文献，不能声称它提供了此处完整Lip证明。
% Fremlin官方已存chap26/47未找到一般coarea全文；265G仅为特殊球面公式。
% Simon作者站NTU2018修订稿前言自称rough lecture draft，仅作背景线索，不增正式书目。
% 本稿的实际自足依赖：13章Besicovitch；15章Hausdorff覆盖量与归一化；
% 16章Borel微分、密度点及McShane；17章近线性分片和加权面积公式；
% 18.1的纯线代余面积恒等式。未使用C1逼近、Sard或19章可整流理论。

\subsection{标量函数}
\label{b1:subsec:180201}

设 \(f:\mathbb R^m\rightarrow\mathbb R\) 为 Lipschitz 函数。对每个实数 \(t\)，水平集 \(f^{-1}(\{t\})\) 是闭集，但它不必是光滑曲面。例如 \(f(x)=|x|\) 在原点不可微，而其正水平集都是球面；一个常值函数则有一个等于整个定义域的水平集。余面积公式要求把这些层按照 \(t\) 的 Lebesgue 测度积分，而不是声称每层都具有同样的几何性质。

\begin{theorem}[标量余面积公式]\label{b1:thm:scalar-coarea}
设 \(f:\mathbb R^m\rightarrow\mathbb R\) 为 Lipschitz 函数，且 \(E\subseteq\mathbb R^m\) Lebesgue 可测。则
\[
 \int_E|\nabla f(x)|\,\mathrm dx
 =
 \int_{\mathbb R}\mathcal H^{m-1}
       \bigl(E\cap f^{-1}(\{t\})\bigr)\,\mathrm dt.
\]
不可微点的梯度约定为零。几乎每个切片在 \(\mathcal H^{m-1}\) 下可测，右端的切片外测度函数为 Lebesgue 可测函数；两端均允许为无穷。
\end{theorem}

这个结论将在下一小节与向量值版本一并证明。当 \(m=1\) 时，\(\mathcal H^0\) 是计数测度，公式已经包含在前章的面积公式中。对于 \(m>1\)，一个新的问题是：Lebesgue 零测集的切片是否仍可忽略？不能仅凭定义域零测，就断言每一层都零测。

下面两个引理同时处理向量值映射。暂设 \(m>n\geq1\)，\(q=m-n\)，且 \(f:\mathbb R^m\rightarrow\mathbb R^n\) 为 \(L\)-Lipschitz 映射。对集合 \(E\)，简记 \(E_y=E\cap f^{-1}(\{y\})\)。在尚未证明切片可测时，\(\mathcal H^q(E_y)\) 表示外测度。

\begin{lemma}\label{b1:lem:null-set-slices}
若 \(N\subseteq\mathbb R^m\) 为 Lebesgue 零测集，则
\[
 \mathcal H^q(N_y)=0
 \quad\text{对几乎每个 }y\in\mathbb R^n.
\]
\end{lemma}
\begin{proof}
记尺度不超过 \(\delta\) 的归一化覆盖量为
\[
 \mathcal H^q_\delta=c_qh^q_\delta,
 \quad c_q=\omega_q/2^q.
\]
由 \(\mathcal H^m(N)=m_m^*(N)=0\)，对每个正整数 \(j\)，可以用闭球 \(B_{j,i}=\overline B(x_{j,i},r_{j,i})\) 覆盖 \(N\)，使
\[
 2r_{j,i}\leq2^{-j},\quad \sum_i r_{j,i}^m<2^{-j}.
\]
从任意直径覆盖得到这样的球覆盖，只需在每个非空覆盖集中选一点，以该集直径为半径作球，并在一开始把尺度及误差预算减小。对直径为零的覆盖集，把半径取为足够小的正数，并使这些额外代价之和仍在预算内。

若切片 \(N_y\) 与球 \(B_{j,i}\) 相交，则 \(y\in f(B_{j,i})\)。因而
\[
 \mathcal H^q_{2^{-j}}(N_y)
 \leq G_j(y):=\omega_q\sum_i r_{j,i}^q
                              \mathbf1_{f(B_{j,i})}(y).
\]
闭球的像紧，故 \(G_j\) 可测。又因
\[
 f(B_{j,i})\subseteq
 \overline B\bigl(f(x_{j,i}),Lr_{j,i}\bigr),
\]
有
\[
 \int_{\mathbb R^n}G_j(y)\,\mathrm dy
 \leq\omega_q\omega_n L^n\sum_i r_{j,i}^{q+n}
 <\omega_q\omega_n L^n2^{-j}.
\]
由 Tonelli 定理，\(\sum_jG_j(y)<\infty\) 几乎处处，因此 \(G_j(y)\to0\) 几乎处处。与此同时，\(\mathcal H^q_{2^{-j}}(N_y)\) 随 \(j\) 增加趋于 \(\mathcal H^q(N_y)\)。结合上述逐点上界，得到所需零测性。
\end{proof}

这个证明没有预先使用切片质量的可测性，而是先构造了可测的上界函数。接着可以确定真正的切片质量何时可测。

\begin{lemma}\label{b1:lem:slice-measure-measurable}
若 \(K\subseteq\mathbb R^m\) 紧，则函数
\[
 y\longmapsto\mathcal H^q(K_y)
\]
为 Borel 可测函数。若 \(E\subseteq\mathbb R^m\) 仅 Lebesgue 可测，则 \(E_y\) 对几乎每个 \(y\) 都是 \(\mathcal H^q\) 可测集，且 \(y\mapsto\mathcal H^q(E_y)\) 为 Lebesgue 可测函数。
\end{lemma}
\begin{proof}
令 \(\mathscr U\) 为全部有理中心、有理半径开球的有限并所成的可数集合族。对固定 \(\delta>0\)，考虑所有有限族 \(U_1,\ldots,U_a\in\mathscr U\)，其中 \(\operatorname{diam}U_i<\delta\)，并以
\[
 c_q\sum_{i=1}^a(\operatorname{diam}U_i)^q
\]
作为覆盖代价。对给定的这样一族，它覆盖 \(K_y\) 的条件是
\[
 y\notin f\left(K\setminus\bigcup_{i=1}^aU_i\right).
\]
括号内的集合紧，其像也紧，故允许的 \(y\) 组成开集。对所有允许的有限覆盖取代价下确界，得到函数 \(F_\delta(y)\)。允许空族覆盖空切片。由于候选族可数，集合 \(\{F_\delta<a\}\) 是可数个开集之并，所以 \(F_\delta\) 为 Borel 函数。

说明其极限确为 Hausdorff 测度。所有候选都是直径小于 \(\delta\) 的覆盖，故
\[
 \mathcal H^q_\delta(K_y)\leq F_\delta(y).
\]
反过来，把 \(K_y\) 的任意直径不超过 \(\delta/2\) 的可数覆盖略微扩大为开集，可以保持每个直径小于 \(\delta\)，并使总代价增加任意小的量；这里使用 \(q>0\) 及幂函数的连续性。紧性给出有限子覆盖。再用有理开球基细化这个有限开覆盖，并把落在同一个原开集内的有限个球合并，就得到 \(\mathscr U\) 中的有限覆盖，代价不增加。因此对每个 \(\eta>0\)，
\[
 F_\delta(y)\leq\mathcal H^q_{\delta/2}(K_y)+\eta.
\]
令 \(\eta\downarrow0\)，再令 \(\delta\downarrow0\)，得到
\[
 \mathcal H^q(K_y)=\lim_{j\to\infty}F_{1/j}(y).
\]
所以紧集情形的切片质量为 Borel 函数。

对于一般 \(E\)，由紧内正则性可取递增紧集 \(K_j\subseteq E\)，使
\[
 E=\bigcup_jK_j\cup N,\quad m_m(N)=0.
\]
前一引理说明 \(N_y\) 几乎处处为 \(\mathcal H^q\) 零测集。对这些 \(y\)，切片 \(E_y\) 是可测集，并且
\[
 \mathcal H^q(E_y)=\lim_{j\to\infty}\mathcal H^q\bigl((K_j)_y\bigr).
\]
于是它在一个 Lebesgue 零测集以外等于可测函数；例外集上的外测度取值即使无穷，也不影响完备 Lebesgue 测度下的可测性。
\end{proof}

\subsection{向量值映射}
\label{b1:subsec:180202}

沿用\cref{b1:def:coarea-jacobian}的记号
\[
 J_nf(x)
 =\sqrt{\det\bigl(Df(x)Df(x)^{\mathsf T}\bigr)}
\]
并在不可微点取零。它非负、Borel 可测，且不超过 \(L^n\)。先处理余面积 Jacobi 行列式消失的部分；与前章不同，这里不要求该部分的像在目标空间中零测。

\begin{lemma}\label{b1:lem:critical-set-slices}
设
\[
 C=\{\,x\in\mathcal D_f:\operatorname{rank}Df(x)<n\,\}.
\]
则 \(\mathcal H^q(C_y)=0\) 对几乎每个 \(y\in\mathbb R^n\) 成立。
\end{lemma}
\begin{proof}
固定正整数 \(R\)，只考虑 \(C_R=C\cap\overline B(0,R)\)。给定 \(x\in C_R\) 及 \(0<\varepsilon\leq1\)，可微性保证充分小的 \(\rho>0\) 满足
\[
 f\bigl(\overline B(x,\rho)\bigr)
 \subseteq f(x)+Df(x)\bigl(\overline B(0,\rho)\bigr)
                   +\overline B(0,\varepsilon\rho).
\]
记 \(r=\operatorname{rank}Df(x)\leq n-1\)。选取目标空间的正交坐标，使前 \(r\) 个方向张成 \(Df(x)\) 的像。由于 \(\|Df(x)\|\leq L\)，右端包含于一个长方体，前 \(r\) 条边的长度不超过 \(2(L+\varepsilon)\rho\)，其余边的长度不超过 \(2\varepsilon\rho\)。因此
\[
 m_n\Bigl(f\bigl(\overline B(x,\rho)\bigr)\Bigr)
 \leq 2^n(1+L)^{n-1}\varepsilon\rho^n.
\]
闭球的像紧，故这里确实是在估计可测集的体积。

对每个 \(j\)，令 \(\varepsilon_j=2^{-j}\)。为 \(C_R\) 中每一点选一个使上述估计成立、且 \(2\rho\leq2^{-j}\) 的中心闭球。由\cref{b1:thm:besicovitch-covering}，可从中选出可数个闭球 \(B_{j,i}=\overline B(x_{j,i},r_{j,i})\) 覆盖 \(C_R\)，并将它们分为至多 \(5^m\) 个两两不交的族。球均包含于 \(B(0,R+1)\)，所以
\[
 \sum_i r_{j,i}^m\leq5^m(R+1)^m.
\]
如前一小节那样，构造
\[
 H_j(y)=\omega_q\sum_i r_{j,i}^q
                                \mathbf1_{f(B_{j,i})}(y).
\]
有逐点上界 \(\mathcal H^q_{2^{-j}}\bigl((C_R)_y\bigr)\leq H_j(y)\)，而上述像体积估计给出
\[
 \int_{\mathbb R^n}H_j(y)\,\mathrm dy
 \leq
 \omega_q2^n(1+L)^{n-1}5^m(R+1)^m\,2^{-j}.
\]
于是 \(\sum_jH_j<\infty\) 几乎处处。令 \(j\to\infty\)，得到 \(\mathcal H^q\bigl((C_R)_y\bigr)=0\) 几乎处处。最后对整数 \(R\) 取可数并，得到结论。
\end{proof}

证明只排除了几乎每层中的正 \(\mathcal H^q\) 测度贡献，并未排除临界纤维中仍有点。因此不能把它改述成“几乎每个水平集都不经过临界点”。

\begin{theorem}[余面积公式]\label{b1:thm:coarea-formula}
设 \(m\geq n\geq1\)，\(f:\mathbb R^m\rightarrow\mathbb R^n\) 为 Lipschitz 映射，且 \(E\subseteq\mathbb R^m\) Lebesgue 可测。则
\[
 \int_EJ_nf(x)\,\mathrm dx
 =
 \int_{\mathbb R^n}
       \mathcal H^{m-n}\bigl(E\cap f^{-1}(\{y\})\bigr)\,\mathrm dy.
\]
右端的切片可测性与质量函数可测性按前述意义成立，两边允许为无穷。
\end{theorem}

\begin{proof}[\cref{b1:thm:coarea-formula}及其标量版本的证明]
当 \(m=n\) 时，\(J_nf=|\det Df|\)，而切片的 \(\mathcal H^0\) 测度为重数，结论就是\cref{b1:thm:area-formula}。以下设 \(q=m-n>0\)。

令 \(G=\{x\in\mathcal D_f:\operatorname{rank}Df(x)=n\}\)。在 \(G\) 外，余面积 Jacobi 行列式为零；不可微点由\cref{b1:lem:null-set-slices}处理，秩亏点由\cref{b1:lem:critical-set-slices}处理。因此只需证明 \(E\cap G\) 上的公式。

在每个 \(x\in G\)，可以选取一个由 \(q\) 个坐标组成的投影 \(P:\mathbb R^m\rightarrow\mathbb R^q\)，使
\[
 B(x)=\begin{pmatrix}P\\ Df(x)\end{pmatrix}
\]
可逆：在 \(Df(x)\) 中选一个非零的 \(n\) 阶坐标子式，再用其余坐标补齐即可。这样的投影只有有限种。对每一种 \(P\)，定义 Lipschitz 映射
\[
 \Phi_P(x)=(Px,f(x))\in\mathbb R^q\times\mathbb R^n.
\]
在其微分可逆的集合上，应用\cref{b1:lem:linear-approximation-pieces}，并把误差取得小于对应线性映射最小伸缩的一半。由近线性的距离比较，可得到可数个 Borel 片，使 \(\Phi_P\) 在每片上为双 Lipschitz 映射。合并有限种 \(P\) 的片，再依次相减，就把 \(G\) 分为两两不交的 Borel 集 \(A_j\)，每片各有这样一个投影 \(P_j\)。

固定一片，并简写 \(A=E\cap A_j\)、\(P=P_j\)、\(\Phi=\Phi_P\)。记
\[
 Z=\Phi(A)\subseteq\mathbb R^q\times\mathbb R^n.
\]
它是 Lebesgue 可测集，因为 \(\Phi\) 是 Lipschitz 映射，且定义域、目标空间同维。片上的逆映射 \(H:Z\rightarrow A\) 为 Lipschitz 映射。逐坐标延拓，得到全空间的 Lipschitz 映射 \(\widetilde H:\mathbb R^m\rightarrow\mathbb R^m\)。

对几乎每个 \(u=(z,y)\in Z\)，\(u\) 是 \(Z\) 的密度点，且 \(\widetilde H\) 在 \(u\) 处可微。置 \(x=H(u)\)，由于 \(x\in A_j\subseteq\mathcal D_f\)，映射 \(\Phi\) 在 \(x\) 处可微。又因
\[
 \Phi\bigl(\widetilde H(v)\bigr)=v\quad(v\in Z),
\]
在密度点处比较双方的微分可得
\[
 D\Phi(x)D\widetilde H(u)=I_m,\quad
 D\widetilde H(u)=B(x)^{-1}.
\]
这里的比较只需密度点的线性唯一性：若一个在 \(u\) 可微的函数在密度一集合上等于零，其微分也为零；否则沿微分不为零的锥形方向会出现一块固定正比例的坏点区域。链式展开则直接来自两映射在相应点的通常可微性。

由\cref{b1:prop:coarea-coordinate-identity}的纯线代形式，上式给出
\[
 J_q\bigl(D_z\widetilde H(z,y)\bigr)
 =
 \frac{J_nf\bigl(H(z,y)\bigr)}
      {|\det B\bigl(H(z,y)\bigr)|}
 \quad\text{在 }Z\text{ 中几乎处处}.
\]
左端使用微分矩阵的前 \(q\) 列。不可微点上补零使这个矩阵函数可测，不改变上述几乎处处等式。

Fubini 定理保证：对几乎每个固定的 \(y\)，截面
\[
 Z_y=\{\,z:(z,y)\in Z\,\}
\]
Lebesgue 可测，并且前述微分关系对几乎每个 \(z\in Z_y\) 成立。映射 \(z\mapsto\widetilde H(z,y)\) 是 \(\mathbb R^q\) 上的 Lipschitz 映射，在 \(Z_y\) 上单射，而且它的像恰好为 \(A\cap f^{-1}(\{y\})\)。在联合可微点处，该映射的微分正是 \(D_z\widetilde H\)。因此由单射面积公式，
\[
 \mathcal H^q\bigl(A\cap f^{-1}(\{y\})\bigr)
 =
 \int_{Z_y}J_q\bigl(D_z\widetilde H(z,y)\bigr)\,\mathrm dz
\]
对几乎每个 \(y\) 成立。

另一方面，对单射映射 \(\Phi|_A\) 使用加权面积公式，权取为 \(J_nf/|\det B|\)，得到
\[
 \int_AJ_nf(x)\,\mathrm dx
 =
 \int_Z
 \frac{J_nf\bigl(H(u)\bigr)}
      {|\det B\bigl(H(u)\bigr)|}\,\mathrm du.
\]
分母在 \(A\) 上非零；涉及的权函数均可测。把矩阵恒等式及切片面积式代入，再用 Tonelli 定理，便得
\[
 \int_AJ_nf(x)\,\mathrm dx
 =
 \int_{\mathbb R^n}\mathcal H^q
                 \bigl(A\cap f^{-1}(\{y\})\bigr)\,\mathrm dy.
\]
最后对两两不交的 \(E\cap A_j\) 求和。几乎每个纤维中的这些部分都可测，因而其 Hausdorff 测度可数相加。再补回已经证明几乎每层零测的 \(E\setminus G\)，得到一般公式。

若 \(n=1\)，微分矩阵只有一行，因此 \(J_1f=|\nabla f|\)。这同时证明了\cref{b1:thm:scalar-coarea}。
\end{proof}

证明没有要求片上的坐标变换在开邻域内可逆。它只用双 Lipschitz 片上的逆映射，并在片像的密度点处识别其微分。这是通常光滑坐标证明在 Lipschitz 情形中需要补充的步骤。

\subsection{可积函数版本}
\label{b1:subsec:180203}

先说明如何限制一个 Lebesgue 可测函数到纤维。设 \(m>n\)。若 \(g\) 有限值且 Lebesgue 可测，可以选一个几乎处处相等的 Borel 代表 \(g_0\)。二者不同的点组成零测集，其切片由\cref{b1:lem:null-set-slices}几乎处处为 Hausdorff 零测集。因此几乎每层上的 \(g\) 都是完备 \(\mathcal H^q\) 可测的，并与 \(g_0\) 给出相同积分。对非负扩展实值函数，以可数个可测水平集作同样处理即可。当 \(m=n\) 时，零测集的像为零测集，所以几乎每个纤维根本不遇到改值点，同样得到结论。

\begin{theorem}[加权余面积公式]\label{b1:thm:weighted-coarea}
在\cref{b1:thm:coarea-formula}的条件下，设 \(g:E\rightarrow[0,\infty]\) Lebesgue 可测。则几乎每个纤维上都有定义良好的非负积分，并且
\[
 \int_Eg(x)J_nf(x)\,\mathrm dx
 =
 \int_{\mathbb R^n}
   \left(\int_{E\cap f^{-1}(\{y\})}g(x)\,
                     \mathrm d\mathcal H^{m-n}(x)\right)\,\mathrm dy.
\]
内积分在例外零测集上任意补值后为 Lebesgue 可测函数，等式允许两端为无穷，并约定 \(0\cdot\infty=0\)。
\end{theorem}
\begin{proof}
若 \(g\) 是非负简单函数，对它的各个可测水平集应用余面积公式，再线性相加即可。对一般 \(g\)，取非负简单函数列 \(g_j\uparrow g\)。上述切片可测性允许删去一个公共零测集，使每个 \(g_j\) 及 \(g\) 都能在余下纤维上积分。逐层应用单调收敛定理，再对外层及定义域上的积分使用单调收敛定理，就得到结论。当 \(m=n\) 时，这也是前章加权面积公式的计数形式。
\end{proof}

若 \(g\) 实值或复值，并满足
\[
 \int_E|g(x)|J_nf(x)\,\mathrm dx<\infty,
\]
先对 \(|g|\) 应用定理，得到几乎每层上的绝对可积性。再对实部、虚部的正负部分分别应用非负公式，便得到同样的等式。若不具有这个条件，不能仅凭某一种迭代积分存在，就交换变号的积分。

这些结论也适用于开集上的局部 Lipschitz 映射。取可数个具有有限 Lipschitz 常数的开球，并把定义域分成从属于它们的不交可测片。在每个球上延拓至全空间后应用公式，最后逐片相加。原映射与延拓在球内局部相同，故微分及余面积 Jacobi 因子一致；补回局部异常集时仍只涉及可数个零测集。

\subsection{几何解释}
\label{b1:subsec:180204}

对投影 \(f(x_1,\ldots,x_m)=(x_1,\ldots,x_n)\)，有 \(J_nf=1\)，纤维是互相平行的 \((m-n)\) 维平面，公式退化为 Fubini 定理。对满秩线性映射，\cref{b1:thm:linear-coarea-formula}说明余面积 Jacobi 因子补偿了法向方向上的伸缩。一般 Lipschitz 映射则在几乎处处使用相应微分矩阵进行同样的补偿。

若 \(n=1\) 且 \(\nabla f(x)\ne0\)，微分核为与 \(\nabla f(x)\) 正交的超平面，\(|\nabla f(x)|\) 测量函数值沿法向变化的速率。沿函数值 \(t\) 积分时，变化快的地方需要更大的补偿因子，才能恢复原空间中的体积。这个解释只涉及正规点的局部微分，不把所有水平集都说成光滑超曲面。

对几乎每个 \(y\)，上面的证明实际上把水平集分成可数个由 \(\mathbb R^{m-n}\) 的子集作 Lipschitz 参数化的部分，再加一个 \(\mathcal H^{m-n}\) 零测集。这个结构是在证明中由满秩坐标片得到的，不是预先假定的水平集正则性。后面的可整流性章节会把这种可数参数化单独命名，并研究它与切空间的关系。

若 \(m_m(E)<\infty\)，则
\[
 \int_{\mathbb R^n}\mathcal H^{m-n}(E_y)\,\mathrm dy
 \leq L^n m_m(E).
\]
因此几乎每个切片的 \((m-n)\) 维测度有限，并且对 \(\lambda>0\)，
\[
 m_n\{\,y:\mathcal H^{m-n}(E_y)>\lambda\,\}
 \leq\lambda^{-1}\int_EJ_nf.
\]
这是一条关于大多数层的估计，而不是每层的统一上界。它将用于下一节的层析积分与水平集估计。

% !TeX root = ../../Book1.tex
\section{应用}
\label{b1:sec:1803}

余面积公式把一个高维积分拆成水平集上的积分，但实际使用时仍须回答三个问题：水平集上该用哪一维测度，是否需要除以梯度，以及梯度为零的参数如何处理。本节从坐标投影和径向分层开始，再讨论水平集面积的平均控制与分布函数。最后把这一方法和第七章已有的 Cavalieri 原理放在一起，说明两种“分层”并不使用同一个积分对象。

% 来源：本书1802的完整余面积公式及1704球面积；主代理实际读取
% Fremlin2016Measure chap26.tex 265G--H的球壳积分与球面积推导。
% 分布推前按本书RN/积分微分工具自行组织；不把临界部分的推前自动宣称为奇异测度。
% 振荡三角波反例为自拟，展示固定Lipschitz常数也不控制全部水平集的统一面积。

\subsection{层析积分}
\label{b1:subsec:180301}

先取坐标投影 \(p:\mathbb R^d\to\mathbb R\)，\(p(x)=x_1\)。它的余面积 Jacobi 因子为 \(|\nabla p|=1\)，水平集是仿射平面 \(\{x_1=t\}\)。该平面上的 \(\mathcal H^{d-1}\) 与其余坐标的 Lebesgue 测度一致，因此对非负可测 \(g\)，余面积公式给出
\[
 \int_{\mathbb R^d}g(x)\,\mathrm dx
 =\int_{\mathbb R}\int_{\mathbb R^{d-1}}
        g(t,z)\,\mathrm dz\,\mathrm dt.
\]
这里得到的就是 Tonelli 公式。若换成任意单位向量方向上的正交投影，只需作一次正交坐标变换。一般余面积公式允许这些直平面随位置弯曲，同时用 Jacobi 因子补偿层与层之间的参数间距。

\begin{proposition}\label{b1:prop:radial-shell-integration}
设 \(d\geq2\)，\(a\in\mathbb R^d\)。对每个非负 Lebesgue 可测函数 \(g\)，
\[
 \int_{\mathbb R^d}g(x)\,\mathrm dx
 =\int_0^\infty\int_{\partial B(a,r)}g(x)\,
       \mathrm d\mathcal H^{d-1}(x)\,\mathrm dr.
\]
若 \(g\) 实值或复值且可积，同一等式成立，内层积分在几乎每个 \(r\) 上绝对可积。
\end{proposition}
\begin{proof}
令 \(u(x)=|x-a|\)。它是 \(1\)-Lipschitz 函数，并且在 \(x\ne a\) 时
\[
 \nabla u(x)=\frac{x-a}{|x-a|},\quad |\nabla u(x)|=1.
\]
单点 \(\{a\}\) 的 \(d\) 维测度为零。在\cref{b1:thm:weighted-coarea}中使用映射 \(u\) 及权 \(g\)，注意正水平集恰为 \(\partial B(a,r)\)，便得到非负版本。对 \(|g|\) 先应用该式，再分实部、虚部的正负部分，就得到可积版本。
\end{proof}

结合\cref{b1:prop:sphere-hausdorff-area}，若 \(g(x)=\varphi(|x-a|)\)，便有
\begin{equation}\label{b1:eq:radial-one-dimensional-integral}
 \int_{\mathbb R^d}\varphi(|x-a|)\,\mathrm dx
 =d\omega_d\int_0^\infty\varphi(r)r^{d-1}\,\mathrm dr.
\end{equation}
公式中的 \(r^{d-1}\) 是半径为 \(r\) 的球面相对于单位球面的面积因子；它不是把球面上的点按 \(d\) 维体积重新计量。\(d=1\) 时，正半径的水平集由两个点组成，\(\mathcal H^0\) 是计数测度，相应公式就是在中心两侧分别积分。

球壳积分还有一个局部读法。若 \(g\in L^1_{\mathrm{loc}}(\mathbb R^d)\)，令
\[
 W(r)=\int_{B(a,r)}g(x)\,\mathrm dx.
\]
对任意有限半径区间，球壳公式表明 \(W\) 是一个可积函数的不定积分。因此在几乎每个 \(r>0\)，
\[
 W'(r)=\int_{\partial B(a,r)}g\,\mathrm d\mathcal H^{d-1}.
\]
“几乎每个”不能在一般可积 \(g\) 下删去。此处控制的是球面积分关于半径的可积性，而不是每张球面上的点态连续性。Fremlin 的《Measure Theory》\cite[265G]{Fremlin2016Measure}直接从参数化推导球壳积分；这里则将它看成余面积公式的径向特例。

\subsection{水平集的 Hausdorff 测度}
\label{b1:subsec:180302}

设 \(U\subseteq\mathbb R^d\) 开，\(u:U\to\mathbb R\) 为 \(L\)-Lipschitz 函数，\(E\subseteq U\) Lebesgue 可测且 \(m_d(E)<\infty\)。令
\[
 A_E(t)=\mathcal H^{d-1}\bigl(E\cap u^{-1}(\{t\})\bigr).
\]
上一节保证切片在几乎每个参数值上可测，并使 \(A_E\) 具有 Lebesgue 可测的版本。余面积公式给出
\begin{equation}\label{b1:eq:level-area-average}
 \int_{\mathbb R}A_E(t)\,\mathrm dt
 =\int_E|\nabla u(x)|\,\mathrm dx
 \leq Lm_d(E).
\end{equation}
于是几乎每个水平集与 \(E\) 的交具有有限 \((d-1)\) 维测度，而且对 \(M>0\)，
\[
 m_1\{t:A_E(t)>M\}\leq\frac{Lm_d(E)}M.
\]
后一个估计来自非负积分的简单下界 \(\int A_E\geq M m_1\{A_E>M\}\)，它控制大面积水平集出现的参数范围。

这并不控制每个水平集。若 \(u\) 在一个开球上恒为 \(c\)，则 \(u^{-1}(\{c\})\) 含有该开球，它的 \((d-1)\) 维测度为无穷；异常参数只有一个，仍与\eqref{b1:eq:level-area-average}相容。即使排除这种平台，固定定义域与 Lipschitz 常数也不能给所有正常水平集一个统一的面积上界。

为看清后一件事，在单位开立方体 \(Q=(0,1)^d\) 上取
\[
 u_N(x)=\frac1N\operatorname{dist}(Nx_1,\mathbb Z),
 \quad N=1,2,ldots.
\]
距离函数为 \(1\)-Lipschitz，所以每个 \(u_N\) 也为 \(1\)-Lipschitz。对 \(0<t<1/(2N)\)，方程 \(u_N(x)=t\) 在第一坐标上恰有 \(2N\) 个解，其余坐标自由。因此水平集由 \(2N\) 个互不相交的单位截面组成，
\[
 A_Q(t)=2N\quad\bigl(0<t<1/(2N)\bigr).
\]
每张这样的水平集面积随 \(N\) 增长，但对应参数区间缩短；其面积积分始终为 \(1\)，也正好等于 \(\int_Q|\nabla u_N|\)。余面积公式描述的是这两种变化之间的平衡。

对于 \(f:U\subseteq\mathbb R^m\to\mathbb R^n\)、\(m\geq n\)，同一论证给出
\[
 \int_{\mathbb R^n}\mathcal H^{m-n}
       \bigl(E\cap f^{-1}(\{y\})\bigr)\,\mathrm dy
 =\int_EJ_nf\leq L^n m_m(E).
\]
在开集上仅局部 Lipschitz 时，可先限制于较小的相对紧参数块。这里仍只得到几乎处处的测度结论；水平集是不是流形，或能否由可数个 Lipschitz 片覆盖，是另一个结构问题。

\subsection{分布函数}
\label{b1:subsec:180303}

现在考虑一个常见的计算目标：已知空间中的密度 \(g(x)\)，希望求数值 \(u(x)\) 的分布。余面积公式的左端带有 \(|\nabla u|\)，而分布的定义直接对 \(g(x)\,\mathrm dx\) 积分，因此需要仔细处理除以梯度的步骤。

设 \(u:U\subseteq\mathbb R^d\to\mathbb R\) 局部 Lipschitz，\(g\geq0\) Lebesgue 可测且 \(\int_Ug<\infty\)。记 \(D_u\) 为真正可微点集，置
\[
 R=\{x\in D_u:|\nabla u(x)|>0\},\quad C=U\setminus R.
\]
在 \(C\) 中，真正临界点与不可微点仍须区分；后者为 Lebesgue 零测集，对这里的 \(g\,\mathrm dx\) 没有质量。定义有限推前测度
\[
 \nu(B)=\int_{u^{-1}(B)}g(x)\,\mathrm dx
 \quad(B\subseteq\mathbb R\text{ Borel}).
\]

\begin{proposition}\label{b1:prop:distribution-coarea-decomposition}
上述推前测度可以写成
\[
 \nu=\rho(t)\,\mathrm dt+\nu_C,
\]
其中
\[
 \rho(t)=\int_{R\cap u^{-1}(\{t\})}
          \frac{g(x)}{|\nabla u(x)|}\,
          \mathrm d\mathcal H^{d-1}(x),
 \quad
 \nu_C(B)=\int_{C\cap u^{-1}(B)}g(x)\,\mathrm dx.
\]
\(\rho\) 在几乎处处定义，并有一个非负可测版本，满足
\(\int_{\mathbb R}\rho=\int_Rg\)。若 \(\int_Cg=0\)，则尾分布函数
\[
 F(t)=\int_{\{x\in U:u(x)>t\}}g(x)\,\mathrm dx
\]
绝对连续，且 \(F'(t)=-\rho(t)\) 几乎处处。
\end{proposition}
\begin{proof}
给定非负 Borel 测试函数 \(\psi\)，在余面积公式中取权
\[
 h(x)=
 \begin{cases}
  \psi(u(x))g(x)/|\nabla u(x)|,&x\in R,\\
  0,&x\in C.
 \end{cases}
\]
非负版本不要求 \(h\) 本身对体积可积。它的乘积 \(h|\nabla u|\) 等于 \(\mathbf1_R\psi(u)g\)，所以
\[
 \int_R\psi(u(x))g(x)\,\mathrm dx
 =\int_{\mathbb R}\psi(t)\rho(t)\,\mathrm dt.
\]
令 \(\psi=1\) 得到可积性，令 \(\psi=\mathbf1_B\) 得到 \(R\) 部分的推前密度；再加上 \(C\) 部分即得测度分解。若 \(\int_Cg=0\)，则
\[
 F(t)=\int_t^\infty\rho(s)\,\mathrm ds.
\]
因 \(\rho\in L^1(\mathbb R)\)，绝对连续性及导数结论由第十六章的积分表示得到。
\end{proof}

这里 \(\nu_C\) 是临界部分的推前记号，不是未经证明就得到的 Lebesgue 分解中的奇异部分。余面积公式在 \(C\) 上的左端为零，它只说明几乎每张纤维中相应的低维测度为零，不能单凭这一点恢复 \(g\,\mathrm dx\) 在 \(C\) 上的全部质量。要进一步说明 \(\nu_C\) 的性质，必须利用具体映射的结构。

平台是一个明确的例子：若 \(u=c\) 在可测集 \(A\) 上成立，且 \(\int_Ag>0\)，则 \(\nu\) 在 \(c\) 至少有这部分原子质量。此时尾分布在 \(c\) 跳跃，不可能由某个普通密度函数完全描述。另一方面，临界推前也可能是无原子的奇异测度，本章习题会由上一章的 Lipschitz 同胚给出实例。

作为没有临界质量的计算例子，在 \(B(0,1)\subseteq\mathbb R^d\) 上取均匀概率密度 \(g=1/\omega_d\)，令 \(u(x)=|x|^2\)。除原点外，\(|\nabla u(x)|=2|x|\)。对 \(0<t<1\)，水平集为半径 \(\sqrt t\) 的球面，故
\[
 \rho(t)=\frac{d\omega_d t^{(d-1)/2}}{2\omega_d\sqrt t}
         =\frac d2t^{d/2-1}.
\]
在区间外密度为零。因此 \(\nu((0,t])=t^{d/2}\)，与半径 \(\sqrt t\) 球的体积比例相同。分母 \(2\sqrt t\) 不可遗漏；水平集面积本身还不是数值 \(|x|^2\) 的概率密度。

\subsection{Cavalieri 原理}
\label{b1:subsec:180304}

第七章的\term*[cavalieri-yuanli]{Cavalieri 原理}[Cavalieri's principle]说明，同方向截面的面积可以积成体积。对非负可测函数 \(v\)，其下方区域
\[
 S_v=\{(x,t):0<t<v(x)\}
\]
从两个坐标方向切分，便给出已经证明的水平集积分公式
\[
 \int v(x)\,\mathrm dx
 =\int_0^\infty m_d\{x:v(x)>t\}\,\mathrm dt.
\]
这里内层测量的是超水平集的 \(d\) 维体积，不是等值集的 \((d-1)\) 维面积；也不需要 \(v\) 连续或可微。严格依据是\cref{b1:prop:layer-cake-formula}，不应把余面积公式的梯度因子移入这个等式。

若改为用等值集分层，在上一小节的记号下，对 \(a<b\) 则有
\begin{equation}\label{b1:eq:slab-volume-coarea}
 \int_{\{a<u\leq b\}}g(x)\,\mathrm dx
 =\int_a^b\rho(t)\,\mathrm dt+\nu_C((a,b]).
\end{equation}
这个式子来自推前测度作用于区间，因而端点质量也处理得清楚。若 \(\nu_C=0\)，薄层的质量由水平集上的 \(g/|\nabla u|\) 积分得到；若有临界质量，则还必须保留最后一项。

例如令 \(u(x)=ax_1\)、\(a>0\)，在柱体 \((0,1)\times D\) 中取 \(g=1\)，其中 \(D\subseteq\mathbb R^{d-1}\) 的体积为 \(M<\infty\)。每个内部水平截面的面积都是 \(M\)，但其数值参数 \(t\) 在 \((0,a)\) 中变化，分布密度为 \(M/a\)。若只积截面面积，将得到 \(aM\) 而不是柱体体积 \(M\)。除以梯度 \(a\)，正是在修正数值参数与空间厚度的比例。

把这一区别记住，比背诵两个形式相似的公式更有用。超水平集分层直接计算函数值下方的体积；余面积分层则在原空间内部沿等值集切开。后者有赖于映射的微分结构，给出更细的几何信息，也带来临界点与几乎处处切片的额外条件。下一章会进一步考察这些低维集合本身，说明何时它们能在忽略零测集后由可数个参数片描述。

% 本章小结（不计入编号）
\begin{chaptersummary}
水平集的预期维数由微分的核给出；当映射为 \(C^1\) 且微分满秩时，局部求逆把这个线性预期实现为真实的图像坐标。余面积因子 \(J_nf=\sqrt{\det(DfDf^{\mathsf T})}\) 则记录法向的体积伸缩，与前章对列向量构造的面积因子相区别。

一般 Lipschitz 证明先用小球覆盖控制零测集与临界集的切片，再在满秩双 Lipschitz 坐标片上依次使用全维换元、纤维面积公式及 Tonelli 定理。切片的可测性也在这一过程中单独建立，并没有把迭代积分预先当作有意义。

球壳公式、切片面积估计与分布密度展示了同一方法的不同用途。除以 Jacobi 因子只恢复非退化部分的体积，临界部分须另行保留；超水平集的 Cavalieri 分层则不涉及这一除法。下一章将把余面积证明中出现的可数参数片抽出来，研究可整流性与近似切空间。
\end{chaptersummary}


\BookExercises

% 第十八章习题临时稿，不另设 BookExercises 标题。
% 八题均为依据本章及前文工具组织的自拟题；不冒认特定教材题号。
% 不使用第十九章的可整流集理论或尚未建立的曲面余面积公式。

\begin{exercise}\label{b1:ex:18-oblique-ball-fibers}
令 \(B=\{\,x\in\mathbb R^3\mid |x|\le1\,\}\)，并定义
\[
 F(x_1,x_2,x_3)=(x_1+x_2,x_3).
\]
\begin{enumerate}
\item\label{b1:item:18-oblique-fiber-length}
计算 \(J_2F\)，写出每条纤维 \(B\cap F^{-1}(a,b)\) 的等距参数式及长度。
\item\label{b1:item:18-oblique-pushforward}
求 \(F_\#(m_3|_B)\) 关于 \(m_2\) 的密度，并直接积分核对其总质量。若将 \(m_3|_B\) 换成球上的均匀概率测度，密度怎样改变？
\end{enumerate}
\end{exercise}
% 来源：自拟；用斜投影下的实际弦长辨别纤维长度与体积推前密度。
\begin{solution}
矩阵 \(DF\) 的两行为 \((1,1,0)\) 与 \((0,0,1)\)，所以
\[
 DF\,DF^{\mathsf T}=\begin{pmatrix}2&0\\0&1\end{pmatrix},
 \quad J_2F=\sqrt2.
\]
全部原像可写成
\[
 x=\left(\frac a2,\frac a2,b\right)
 +s\frac{(1,-1,0)}{\sqrt2}.
\]
右端两项正交，参数方向是单位向量。因此纤维非空当且仅当 \(a^2/2+b^2\le1\)，此时
\[
 |s|\le\sqrt{1-a^2/2-b^2},
 \quad
 \mathcal H^1\bigl(B\cap F^{-1}(a,b)\bigr)
 =2\sqrt{1-a^2/2-b^2}.
\]
等号边界上的纤维是单点，其一维测度为零。

对非负 Borel 函数 \(h\)，在\cref{b1:thm:weighted-coarea}中取权重 \(h\circ F/\sqrt2\)，得到
\[
 \int_Bh\bigl(F(x)\bigr)\,\mathrm dm_3(x)
 =\int_{\mathbb R^2}h(a,b)\rho(a,b)\,\mathrm da\,\mathrm db,
\]
其中
\[
 \rho(a,b)=
 \begin{cases}
 \sqrt2\sqrt{1-a^2/2-b^2},&a^2/2+b^2<1,\\
 0,&a^2/2+b^2\ge1.
 \end{cases}
\]
密度等于弦长除以 \(\sqrt2\)，而不是弦长本身。令 \(a=\sqrt2 u\)、\(b=v\)，再用二维极坐标，便有
\[
 \int_{\mathbb R^2}\rho\,\mathrm dm_2
 =2\int_{u^2+v^2<1}\sqrt{1-u^2-v^2}\,\mathrm du\,\mathrm dv
 =4\pi\int_0^1r\sqrt{1-r^2}\,\mathrm dr
 =\frac{4\pi}{3}.
\]
这与球体积一致。均匀概率测度的推前密度为 \(3\rho/(4\pi)\)。
\end{solution}

\begin{exercise}\label{b1:ex:18-product-logarithmic-density}
在单位方形 \(Q=(0,1)^2\) 上令 \(u(x,y)=xy\)。用余面积公式而不是先对矩形作逐次积分，求 \(u_\#(m_2|_Q)\) 的密度和分布函数。
进一步计算
\[
 \int_Q(xy)^a\,\mathrm dx\,\mathrm dy,
 \quad a>-1.
\]
\end{exercise}
% 来源：自拟；通过双曲线上的弧长和梯度相消导出乘积分布的对数密度。
\begin{solution}
在 \(Q\) 内梯度 \(\nabla u=(y,x)\) 从不为零。对 \(0<t<1\)，水平集可用
\[
 \gamma_t(x)=(x,t/x),
 \quad t<x<1
\]
单射参数化，且
\[
 |\gamma_t'(x)|=\sqrt{1+t^2/x^4},
 \quad \bigl|\nabla u\bigl(\gamma_t(x)\bigr)\bigr|
 =\sqrt{x^2+t^2/x^2}
 =x\sqrt{1+t^2/x^4}.
\]
因此曲线面积公式给出
\[
 \int_{Q\cap u^{-1}(t)}\frac1{|\nabla u|}\,\mathrm d\mathcal H^1
 =\int_t^1\frac{\mathrm dx}{x}=-\log t.
\]
其他水平集为空。将权重 \(h\circ u/|\nabla u|\) 代入\cref{b1:thm:weighted-coarea}，可得对每个非负 Borel 函数 \(h\)，
\[
 \int_Qh(xy)\,\mathrm dx\,\mathrm dy
 =\int_0^1h(t)(-\log t)\,\mathrm dt.
\]
故密度为 \((-\log t)\mathbf1_{(0,1)}(t)\)。方形面积为 \(1\)，它已经是概率密度。

由 \(t\log t\to0\) 可知分布函数为
\[
 G(t)=
 \begin{cases}
 0,&t\le0,\\
 t(1-\log t),&0<t<1,\\
 1,&t\ge1.
 \end{cases}
\]
最后取 \(h(t)=t^a\)。分部积分及 \(t^{a+1}\log t\to0\) 给出
\[
 \int_0^1t^a(-\log t)\,\mathrm dt
 =\frac1{a+1}\int_0^1t^a\,\mathrm dt
 =\frac1{(a+1)^2}.
\]
密度在零点附近无界并不意味着存在原子；这里每个单点的推前质量仍为零。
\end{solution}

\begin{exercise}\label{b1:ex:18-folded-radial-levels}
在圆盘 \(D=\{\,x\in\mathbb R^2\mid |x|\le2\,\}\) 上令
\[
 u(x)=\bigl(|x|-1\bigr)^2.
\]
\begin{enumerate}
\item\label{b1:item:18-folded-level-density}
求 \(0<t<1\) 时水平集的两个分支及其梯度大小，继而计算 \(u_\#(m_2|_D)\) 的密度。
\item\label{b1:item:18-folded-distribution}
若在 \(D\) 上取均匀概率测度，求 \(u\) 的分布函数。说明临界水平 \(t=0\) 是否产生原子，以及为何两个分支的贡献相加而不相减。
\end{enumerate}
\end{exercise}
% 来源：自拟；把单调径向积分改为两个反向分支，区分临界密度发散与原子。
\begin{solution}
函数 \(r\mapsto(r-1)^2\) 在 \([0,2]\) 上为 \(2\)-Lipschitz 函数，故 \(u\) 在 \(D\) 上也是 \(2\)-Lipschitz 函数。除原点外，
\[
 \nabla u(x)=2(|x|-1)\frac{x}{|x|},
 \quad |\nabla u(x)|=2\bigl||x|-1\bigr|.
\]
对 \(0<t<1\)，方程 \(u(x)=t\) 给出两个互不相交的圆，其半径为
\[
 r_-(t)=1-\sqrt t,
 \quad r_+(t)=1+\sqrt t.
\]
在两个圆上梯度大小都等于 \(2\sqrt t\)。因此水平积分为
\[
 \int_{D\cap u^{-1}(t)}\frac1{|\nabla u|}\,\mathrm d\mathcal H^1
 =\frac{2\pi(1-\sqrt t)+2\pi(1+\sqrt t)}{2\sqrt t}
 =\frac{2\pi}{\sqrt t}.
\]
原点、单位圆和外边界均为二维零测集，可以在它们以外应用加权余面积公式，再补回零质量部分。所得推前密度为
\[
 \rho(t)=\frac{2\pi}{\sqrt t}\mathbf1_{(0,1)}(t),
\]
它的总积分为 \(4\pi=m_2(D)\)。

均匀概率密度因此是 \(1/(2\sqrt t)\)，分布函数为
\[
 G(t)=
 \begin{cases}
 0,&t\le0,\\
 \sqrt t,&0<t<1,\\
 1,&t\ge1.
 \end{cases}
\]
也可用环域检验：当 \(0<t<1\) 时，\(\{u\le t\}\) 的面积为
\[
 \pi\bigl((1+\sqrt t)^2-(1-\sqrt t)^2\bigr)=4\pi\sqrt t.
\]
临界水平的原像是单位圆，二维面积为零，所以 \(t=0\) 没有原子。两支上径向导数的符号相反，但余面积因子是梯度的模；两片原像各自携带正面积，必须相加。
\end{solution}

\begin{exercise}\label{b1:ex:18-two-circle-vector-fibers}
令
\[
 \Omega=\{\,(x,y,z)\in\mathbb R^3\mid x^2+y^2<1,\ |z|<1\,\},
 \quad F(x,y,z)=(x^2+y^2,z^2).
\]
计算 \(J_2F\) 和 \(0<s,t<1\) 时的纤维，求下列两个有限测度关于 \(m_2\) 的密度：
\[
 F_\#(m_3|_\Omega),
 \quad F_\#\bigl(|z|m_3|_\Omega\bigr).
\]
说明临界集在推前中是否留下额外质量。
\end{exercise}
% 来源：自拟；以两个圆组成的纤维练习向量值余面积和权重，显式处理退化轴。
\begin{solution}
记 \(r=\sqrt{x^2+y^2}\)。微分矩阵及其行 Gram 矩阵分别为
\[
 DF=\begin{pmatrix}2x&2y&0\\0&0&2z\end{pmatrix},
 \quad DF\,DF^{\mathsf T}=\begin{pmatrix}4r^2&0\\0&4z^2\end{pmatrix},
\]
故 \(J_2F=4r|z|\)。该映射在有界凸圆柱上 Lipschitz，退化集是轴 \(r=0\) 与平面 \(z=0\) 在 \(\Omega\) 内的部分，两者都是三维零测集。

对 \(0<s,t<1\)，纤维由位于 \(z=\sqrt t\) 与 \(z=-\sqrt t\) 的两个圆组成，半径均为 \(\sqrt s\)，总长度为 \(4\pi\sqrt s\)。纤维上 \(J_2F=4\sqrt s\sqrt t\)，所以
\[
 \int_{\Omega\cap F^{-1}(s,t)}\frac1{J_2F}\,\mathrm d\mathcal H^1
 =\frac{\pi}{\sqrt t},
 \quad
 \int_{\Omega\cap F^{-1}(s,t)}\frac{|z|}{J_2F}\,\mathrm d\mathcal H^1
 =\pi.
\]
先在非退化集上用\cref{b1:thm:weighted-coarea}，再补回零质量临界集，便得两个密度
\[
 \rho_0(s,t)=\frac{\pi}{\sqrt t}\mathbf1_{(0,1)^2}(s,t),
 \quad \rho_1(s,t)=\pi\mathbf1_{(0,1)^2}(s,t).
\]
第一项的总积分为 \(2\pi\)，等于圆柱体积；第二项的总积分为 \(\pi\)，等于 \(\pi\int_{-1}^1|z|\,\mathrm dz\)。临界集的原测度为零，故它不产生附加推前质量。两个圆的贡献都已计入；只保留 \(z>0\) 的圆会把两个答案都减半。
\end{solution}

\begin{exercise}\label{b1:ex:18-distance-to-segment}
设 \(L>0\)，\(S=[0,L]\times\{0\}\subset\mathbb R^2\)，令 \(d(x)=\operatorname{dist}(x,S)\)。对 \(R>0\)，记 \(S_R=\{\,x\mid d(x)<R\,\}\)。
\begin{enumerate}
\item\label{b1:item:18-segment-level-length}
逐段描述 \(d^{-1}(t)\)，计算其长度，并用余面积公式求 \(m_2(S_R)\)。
\item\label{b1:item:18-segment-weighted-layer}
对非负 Borel 函数 \(h:(0,R)\rightarrow[0,\infty]\)，求 \(\int_{S_R\setminus S}h\bigl(d(x)\bigr)\,\mathrm dm_2(x)\)。特别地，求 \(\int_{S_R}d(x)\,\mathrm dm_2(x)\)，并核对 \(m_2(S_R)\) 对 \(R\) 的导数与边界长度的关系。
\end{enumerate}
\end{exercise}
% 来源：自拟；以非球形平行集计算水平长度、距离矩和薄层体积，不调用周长或 BV 理论。
\begin{solution}
写 \(x=(a,b)\)。由到线段的最近点位置，有
\[
 d(a,b)=
 \begin{cases}
 \sqrt{a^2+b^2},&a<0,\\
 |b|,&0\le a\le L,\\
 \sqrt{(a-L)^2+b^2},&a>L.
 \end{cases}
\]
距离函数为 \(1\)-Lipschitz 函数，由这些公式可见 \(|\nabla d|=1\) 在 \(\mathbb R^2\setminus S\) 上几乎处处成立。

对 \(t>0\)，水平集由 \(b=\pm t\)、\(0\le a\le L\) 的两条线段，以及分别以 \((0,0)\)、\((L,0)\) 为圆心的外侧半圆组成。四个连接点不影响长度，故
\[
 \mathcal H^1\bigl(d^{-1}(t)\bigr)=2L+2\pi t.
\]
又 \(m_2(S)=0\)，所以余面积公式给出
\[
 m_2(S_R)=\int_0^R(2L+2\pi t)\,\mathrm dt
 =2LR+\pi R^2.
\]
这对应一个长为 \(L\)、宽为 \(2R\) 的矩形，加上两个半圆。

同样，由加权版本得到
\[
 \int_{S_R\setminus S}h\bigl(d(x)\bigr)\,\mathrm dm_2(x)
 =\int_0^Rh(t)(2L+2\pi t)\,\mathrm dt.
\]
非负性允许两端同时为无穷。取 \(h(t)=t\)，得到
\[
 \int_{S_R}d(x)\,\mathrm dm_2(x)
 =LR^2+\frac{2\pi}{3}R^3.
\]
因为 \(\partial S_R=d^{-1}(R)\)，上述面积表达式还给出
\[
 \frac{\mathrm d}{\mathrm dR}m_2(S_R)
 =2L+2\pi R=\mathcal H^1(\partial S_R).
\]
这里边界的连接点无需拥有统一的曲率；水平集长度已经足以计算薄层体积。
\end{solution}

\begin{exercise}\label{b1:ex:18-cube-oblique-sections}
令 \(C=[0,1]^3\)，\(u(x)=x_1+x_2+x_3\)，并记 \(r_+=\max\{r,0\}\)。证明对每个 \(t\in\mathbb R\)，
\[
 \mathcal H^2\bigl(C\cap u^{-1}(t)\bigr)
 =\frac{\sqrt3}{2}\sum_{j=0}^3(-1)^j\binom3j(t-j)_+^2.
\]
求截面积的最大值和取得它的水平，写出最大截面的顶点。
提示：先求 \(V(t)=m_3(C\cap\{u\le t\})\)，再解释为什么从余面积公式得到的几乎处处等式能够延拓到每个 \(t\)。
\end{exercise}
% 来源：自拟；把容斥体积、余面积公式与切片连续性相连，避免将 a.e. 导数关系直接误写为处处。
\begin{solution}
在正八分体中，单纯形 \(\{x_i\ge0,\ \sum_i x_i\le t\}\) 的体积为 \(t_+^3/6\)。对违反 \(x_i\le1\) 的坐标作容斥；若指定 \(j\) 个坐标超过 \(1\)，将这些坐标各减去 \(1\)，所得体积为 \((t-j)_+^3/6\)。边界平面的三维测度为零，故
\[
 V(t)=\frac16\sum_{j=0}^3(-1)^j\binom3j(t-j)_+^3.
\]
因此 \(V\) 处处可微，且
\[
 V'(t)=\frac12\sum_{j=0}^3(-1)^j\binom3j(t-j)_+^2.
\]
因为 \(|\nabla u|=\sqrt3\)，余面积公式首先说明
\[
 \mathcal H^2\bigl(C\cap u^{-1}(t)\bigr)=\sqrt3V'(t)
\]
对几乎每个 \(t\) 成立。

为补齐其余水平，把截面写成图像
\[
 (a,b)\longmapsto(a,b,t-a-b),
 \quad (a,b)\in D_t,
\]
其中 \(D_t=\{\,(a,b)\in[0,1]^2\mid t-1\le a+b\le t\,\}\)。该参数映射的面积因子为 \(\sqrt3\)，故截面积等于 \(\sqrt3m_2(D_t)\)。若 \(t_n\to t\)，则除 \(a+b=t-1\) 和 \(a+b=t\) 两条零测直线外，\(\mathbf1_{D_{t_n}}\to\mathbf1_{D_t}\)。控制收敛定理给出 \(m_2(D_{t_n})\to m_2(D_t)\)。截面积与上面的 \(\sqrt3V'\) 都连续，而它们几乎处处相等，因而处处相等。

在 \(0\le t\le1\) 上，截面积为 \(\sqrt3t^2/2\)。在 \(1\le t\le2\) 上，它等于
\[
 \sqrt3\left(-t^2+3t-\frac32\right)
 =\sqrt3\Bigl(\frac34-\bigl(t-\frac32\bigr)^2\Bigr).
\]
由立方体关于 \(x\mapsto(1,1,1)-x\) 的对称性，\([2,3]\) 上的情形与 \([0,1]\) 相同。于是最大值为 \(3\sqrt3/4\)，且只在 \(t=3/2\) 取得。该截面的六个顶点是 \((1,1/2,0)\) 的全部排列，相邻顶点的距离为 \(1/\sqrt2\)，组成正六边形。
\end{solution}

\begin{exercise}\label{b1:ex:18-transverse-singularity}
设整数 \(1\le k<n\)，\(\alpha>0\)，并令
\[
 \Omega=B^k(0,1)\times(0,1)^{n-k},
 \quad f(y,z)=|y|^{-\alpha}\quad(y\ne0).
\]
在零测集 \(y=0\) 上赋值为 \(0\)。利用坐标投影的余面积公式，求所有 \(p\ge1\) 中使 \(f\in L^p(\Omega)\) 成立的 \(p\)，并在可积时计算 \(\|f\|_p^p\)。再求
\[
 m_n\{\,x\in\Omega\mid f(x)>\lambda\,\},
 \quad \lambda>0,
\]
并解释为何在 \(q=k/\alpha\) 时，\(\int_\Omega f^q\,\mathrm dm_n=\infty\)，却仍有
\[
 \sup_{\lambda>0}\lambda^q
 m_n\{\,x\in\Omega\mid f(x)>\lambda\,\}<\infty.
\]
\end{exercise}
% 来源：自拟；连接线性纤维积分、奇性横向维数和临界分布尾估计，不引入未定义的函数空间。
\begin{solution}
对投影 \(P(y,z)=y\)，有 \(J_kP=1\)。当 \(y\in B^k(0,1)\) 时，纤维为 \(\{y\}\times(0,1)^{n-k}\)，其 \((n-k)\) 维 Hausdorff 测度为 \(1\)。由加权余面积公式及 \(k\) 维径向积分，
\[
 \int_\Omega f^p\,\mathrm dm_n
 =\int_{B^k(0,1)}|y|^{-\alpha p}\,\mathrm dm_k(y)
 =k\omega_k\int_0^1r^{k-1-\alpha p}\,\mathrm dr.
\]
这里 \(\omega_k=m_k\bigl(B^k(0,1)\bigr)\)。所以对 \(p\ge1\)，可积的充要条件是 \(\alpha p<k\)，且此时
\[
 \|f\|_p^p=\frac{k\omega_k}{k-\alpha p}.
\]
当 \(\alpha p=k\) 时积分按 \(\int_0^1\mathrm dr/r\) 发散；更大的指数也发散。

对 \(\lambda\ge1\)，不等式 \(|y|^{-\alpha}>\lambda\) 等价于 \(|y|<\lambda^{-1/\alpha}\)，除去无关的零测赋值集后，投影的每条纤维仍有测度 \(1\)。因而
\[
 m_n\{f>\lambda\}
 =\begin{cases}
 \omega_k,&0<\lambda<1,\\
 \omega_k\lambda^{-k/\alpha},&\lambda\ge1.
 \end{cases}
\]
取 \(q=k/\alpha\)，上式给出所求上确界恰为 \(\omega_k\)，但前面的积分计算对任意正指数都适用，故 \(\int f^q=\infty\)。这一区别来自临界幂次产生的对数发散。控制可积性的维数是奇异平面的横向维数 \(k\)，而不是环境维数 \(n\)。
\end{solution}

\begin{exercise}\label{b1:ex:18-critical-singular-continuous-pushforward}
取紧且无处稠密的集合 \(K\subset[0,1]\)，满足 \(a=m_1(K)\in(0,1)\)。令
\[
 f(t)=\int_0^t\mathbf1_{[0,1]\setminus K}(s)\,\mathrm ds,
 \quad \ell=1-a,
 \quad u(x,y)=f(x)\quad((x,y)\in Q=[0,1]^2).
\]
可使用\cref{b1:ex:17-singular-lipschitz-homeomorphism}中已经建立的结论：\(f\) 是 \([0,1]\) 到 \([0,\ell]\) 的 \(1\)-Lipschitz 同胚，且 \(m_1\bigl(f(K)\bigr)=0\)。
\begin{enumerate}
\item\label{b1:item:18-critical-measure-decomposition}
证明
\[
 u_\#(m_2|_Q)=m_1|_{[0,\ell]}+\nu,
 \quad \nu=f_\#(m_1|_K),
\]
其中 \(\nu\) 是无原子的奇异测度，总质量为 \(a\)。
\item\label{b1:item:18-critical-coarea-comparison}
计算每个水平集 \(Q\cap u^{-1}(t)\) 的长度，并比较
\[
 \int_Qh(u)|\nabla u|\,\mathrm dm_2,
 \quad \int_Qh(u)\,\mathrm dm_2
\]
对非负 Borel 函数 \(h\) 的表达式。说明为什么在零梯度处规定 \(1/|\nabla u|=0\)，不能把余面积公式变成全部体积的推前密度公式。
\end{enumerate}
\end{exercise}
% 来源：自拟；把第十七章的单射压缩构造用于临界推前，展示奇异部分不必由平台原子组成。
\begin{solution}
由于 \(u\) 不依赖第二坐标，而该坐标区间长度为 \(1\)，Fubini 定理给出
\[
 u_\#(m_2|_Q)=f_\#(m_1|_{[0,1]}).
\]
记 \(E=[0,1]\setminus K\)。在 \(E\) 内除去区间端点后，\(f\) 局部是斜率为 \(1\) 的仿射函数，所以 \(f'=1\)。由一维加权面积公式及 \(f\) 单射性，对每个非负 Borel 函数 \(h\)，
\[
 \int_Eh\bigl(f(x)\bigr)\,\mathrm dx
 =\int_{f(E)}h(t)\,\mathrm dt
 =\int_0^\ell h(t)\,\mathrm dt.
\]
最后一步用了 \([0,\ell]\setminus f(E)=f(K)\) 为零测集。因此 \(f_\#(m_1|_E)=m_1|_{[0,\ell]}\)，分解即得。

测度 \(\nu\) 集中在零测集 \(f(K)\) 上，故关于 Lebesgue 测度奇异，其总质量为 \(m_1(K)=a\)。又因为 \(f\) 单射，每个单点的原像至多是一个单点，因而
\[
 \nu(\{t\})=m_1\bigl(K\cap f^{-1}(\{t\})\bigr)=0.
\]
这证明它没有原子，而不是仅仅说明没有明显的平台区间。

对每个 \(0\le t\le\ell\)，水平集恰为竖直线段
\[
 Q\cap u^{-1}(t)=\{f^{-1}(t)\}\times[0,1],
\]
其长度为 \(1\)；其他水平集为空。另一方面，\(f'=\mathbf1_E\) 几乎处处，故 \(|\nabla u(x,y)|=\mathbf1_E(x)\) 几乎处处。余面积公式与第一问分别给出
\[
 \int_Qh(u)|\nabla u|\,\mathrm dm_2
 =\int_0^\ell h(t)\,\mathrm dt,
\]
以及
\[
 \int_Qh(u)\,\mathrm dm_2
 =\int_0^\ell h(t)\,\mathrm dt+\int h(t)\,\mathrm d\nu(t).
\]
集合 \(K\times[0,1]\) 的二维测度为 \(a\)，而梯度在那里几乎处处为零。若把零梯度处的倒数设为零，就会删掉这部分体积，得到的仅是第一问中的绝对连续部分。临界值集合 \(f(K)\) 虽然是一维零测集，其原像仍能携带正体积，并产生无原子的奇异推前质量。
\end{solution}
